\documentclass[../../main.tex]{subfiles}
\begin{document}
\graphicspath{{figures/}{lectures/03_logistic/figures/}}
\section{Classification: Logistic Regression}
\begin{frame}[plain]
  \vfill\begin{center}{\Huge \textcolor{blue}{Classification: Logistic Regression}}\end{center}\vfill
\end{frame}
% Title Slide



%--------------------------------------------------
\begin{frame}{Classification: Introduction}
\begin{itemize}
    \item In regression, the response variable $Y$ is \textbf{quantitative}.  
    \item In many applications, $Y$ is instead \textbf{qualitative} (categorical).  
    \begin{itemize}
        \item Examples: eye color, disease presence, fraud detection.  
    \end{itemize}
    \item Predicting a qualitative response is called \textbf{classification}.  
    \item A classifier assigns each observation to a category (class).  
    \item Often we estimate probabilities of class membership, and then classify.  
\end{itemize}
\end{frame}

%--------------------------------------------------
\begin{frame}{Classification: Motivation}
Real-world classification problems:
\begin{enumerate}
    \item \textbf{Medicine:} Predicting which disease a patient has from symptoms.  
    \item \textbf{Finance:} Detecting fraudulent transactions from user behavior.  
    \item \textbf{Biology:} Identifying whether DNA mutations are deleterious.  
\end{enumerate}
\begin{itemize}
    \item Training data: $(x_1,y_1), \dots, (x_n,y_n)$ with $y_i \in \{1,\dots,K\}$.  
    \item Goal: Build a classifier that predicts well on unseen test data.  
    \item Methods include: \textbf{logistic regression, LDA, QDA, naive Bayes, KNN}.  
\end{itemize}
\end{frame}



%--------------------------------------------------
\begin{frame}{Why Not Linear Regression? (I)}
\begin{itemize}
    \item Linear regression assumes a \textbf{quantitative} response.  
    \item For classification, $Y$ is \textbf{categorical} (e.g. stroke, drug overdose, epileptic seizure).  
    \item One could encode categories as numbers:
    \[
    Y =
    \begin{cases}
    1 & \text{if stroke}, \\
    2 & \text{if drug overdose}, \\
    3 & \text{if epileptic seizure}.
    \end{cases}
    \]
    \item Problem: This encoding imposes an \textbf{arbitrary ordering} of classes.  
    \item Different codings lead to different regression models and predictions.  
    \item There is generally no meaningful numeric ordering for qualitative outcomes.  
\end{itemize}
\end{frame}

%--------------------------------------------------
\begin{frame}{Why Not Linear Regression? (II)}
\begin{itemize}
    \item For a binary outcome (e.g. $Y=0$ for stroke, $Y=1$ for drug overdose),  
    linear regression can be fit, but issues remain:
    \begin{itemize}
        \item Predictions can fall outside the $[0,1]$ range.  
        \item Estimates are not proper probabilities.  
        \item Interpretation is problematic.  
    \end{itemize}
    \item Summary of issues:
    \begin{enumerate}
        \item Cannot accommodate more than two classes meaningfully.  
        \item Even with two classes, does not yield valid probability estimates.  
    \end{enumerate}
    \item $\Rightarrow$ Need a proper classification model (e.g. logistic regression, LDA).  
\end{itemize}
\end{frame}


%--------------------------------------------------
\begin{frame}{Motivation for Logistic Regression}
\begin{itemize}
    \item Goal: Model $p(X) = \Pr(Y=1|X)$.  
    \item Linear probability model: $p(X) = \beta_0 + \beta_1 X$.
    \begin{itemize}
        \item Issues: predictions $<0$ or $>1$, nonsensical as probabilities.  
    \end{itemize}
    \item We require a function $p(X)$ that:
    \begin{enumerate}
        \item Maps $\mathbb{R} \to (0,1)$.  
        \item Is flexible, but allows interpretable parameters.  
    \end{enumerate}
    \item Solution: the \textbf{logistic function}.
\end{itemize}
\end{frame}

%--------------------------------------------------
\begin{frame}{The Logistic Function}
We want to model the probability $p(x) = \Pr(Y=1 \mid X=x)$.  
\[
p(x) = \frac{1}{1 + e^{-(x-\mu)/s}},
\]
where:
\begin{itemize}
    \item $\mu$ = location parameter (midpoint, where $p(\mu)=1/2$).  
    \item $s$ = scale parameter (controls slope at midpoint).  
\end{itemize}
\[
\lim_{x\to -\infty} p(x) = 0, \qquad \lim_{x\to +\infty} p(x) = 1.
\]
\begin{itemize}
    \item This function maps $\mathbb{R} \to (0,1)$, suitable for probabilities.  
\end{itemize}
\end{frame}

%--------------------------------------------------
\begin{frame}{Reparameterization}
Rewrite logistic function in regression form:
\[
p(x) = \frac{1}{1 + e^{-(\beta_0 + \beta_1 x)}}.
\]
where:
\[
\beta_0 = -\mu/s \quad \text{(intercept)}, \qquad \beta_1 = 1/s \quad \text{(slope)}.
\]
Interpretation:
\begin{itemize}
    \item $\beta_0$ = log-odds when $x=0$.  
    \item $\beta_1$ = change in log-odds per unit increase in $x$.  
\end{itemize}
Thus:
\[
\text{logit}(p(x)) = \log \frac{p(x)}{1-p(x)} = \beta_0 + \beta_1 x.
\]
\end{frame}

%--------------------------------------------------
\begin{frame}{The Logistic Regression Model}
Generalize to multiple predictors $X \in \mathbb{R}^p$:
\[
p(x) = \Pr(Y=1 \mid X=x) = \frac{1}{1 + e^{-(\beta_0 + \beta^T x)}}.
\]
Key properties:
\begin{itemize}
    \item $p(x)$ is always between 0 and 1.  
    \item $\log \dfrac{p(x)}{1-p(x)} = \beta_0 + \beta^T x$ (linear in predictors).  
    \item Parameters $(\beta_0, \beta)$ estimated by maximum likelihood.  
\end{itemize}
This is the foundation of \textbf{classification via logistic regression}.
\end{frame}



%--------------------------------------------------
\begin{frame}{The Logistic Function}
\[
p(X) = \frac{e^{\beta_0 + \beta^T X}}{1 + e^{\beta_0 + \beta^T X}}.
\]
\begin{itemize}
    \item Always $\in (0,1)$.  
    \item Symmetric S-shaped curve.  
    \item $\lim_{z\to -\infty} p(z)=0$, $\lim_{z\to +\infty} p(z)=1$.  
\end{itemize}
\[
\text{Logistic regression model: } \quad Y_i \sim \text{Bernoulli}(p_i), \quad 
p_i = \frac{e^{\beta_0 + \beta^T x_i}}{1 + e^{\beta_0 + \beta^T x_i}}.
\]
\end{frame}

%--------------------------------------------------
\begin{frame}{Odds and Log-Odds}
Define odds of success:
\[
\text{odds}(X) = \frac{p(X)}{1 - p(X)}.
\]
From logistic model:
\[
\frac{p(X)}{1 - p(X)} = e^{\beta_0 + \beta^T X}.
\]
Taking logarithm:
\[
\log\left(\frac{p(X)}{1 - p(X)}\right) = \beta_0 + \beta^T X.
\]
\begin{itemize}
    \item The \textbf{log-odds (logit)} is linear in $X$.  
    \item Interpretation: coefficients act multiplicatively on the odds.  
\end{itemize}
\end{frame}

%--------------------------------------------------
\begin{frame}{Interpretation of Coefficients}
\begin{itemize}
    \item For predictor $X_j$, coefficient $\beta_j$ means:
    \[
    \log\left(\frac{p(X)}{1 - p(X)}\right) \;\; \text{increases by } \beta_j \;\; \text{for a one-unit increase in } X_j.
    \]
    \item Equivalently, odds are multiplied by $e^{\beta_j}$.  
    \item Example: $\beta_j = 0.7 \Rightarrow e^{0.7} \approx 2.01$, doubling the odds per unit increase in $X_j$.  
\end{itemize}
\end{frame}

%--------------------------------------------------
\begin{frame}{Likelihood for Logistic Regression}
Data: $\{(x_i,y_i)\}_{i=1}^N$, with $y_i \in \{0,1\}$.
\[
p_i = \Pr(Y_i=1|x_i) = \frac{e^{\beta_0 + \beta^T x_i}}{1+e^{\beta_0 + \beta^T x_i}}.
\]
Likelihood:
\[
L(\beta) = \prod_{i=1}^N p_i^{y_i} (1-p_i)^{1-y_i}.
\]
Log-likelihood:
\[
\ell(\beta) = \sum_{i=1}^N \Big[ y_i(\beta_0 + \beta^T x_i) - \log(1 + e^{\beta_0 + \beta^T x_i}) \Big].
\]
\end{frame}

%--------------------------------------------------
\begin{frame}{MLE: First-Order Conditions}
Gradient (score equations):
\[
\frac{\partial \ell(\beta)}{\partial \beta} 
= \sum_{i=1}^N (y_i - p_i)x_i.
\]
Setting gradient $=0$:
\[
\sum_{i=1}^N (y_i - p_i)x_i = 0.
\]
\begin{itemize}
    \item No closed form solution.  
    \item Must use iterative methods: Newton–Raphson or IRLS (iteratively reweighted least squares).  
\end{itemize}
\end{frame}

%--------------------------------------------------
\begin{frame}{Newton--Raphson for Maximum Likelihood}
Goal: maximize the log-likelihood $\ell(\beta)$.  
\begin{itemize}
    \item General iterative scheme:
    \[
    \beta^{(t+1)} = \beta^{(t)} - \Big[ H(\beta^{(t)}) \Big]^{-1} \nabla \ell(\beta^{(t)}),
    \]
    where:
    \begin{itemize}
        \item $\nabla \ell(\beta)$ = gradient (score vector).  
        \item $H(\beta)$ = Hessian matrix (second derivatives).  
    \end{itemize}
    \item Intuition:
    \begin{itemize}
        \item Uses a quadratic approximation of $\ell(\beta)$.  
        \item Updates parameters by moving towards stationary point.  
        \item Converges quickly near the optimum.  
    \end{itemize}
\end{itemize}
\end{frame}

%--------------------------------------------------
\begin{frame}{Newton--Raphson for Logistic Regression}
For logistic regression:
\[
\nabla \ell(\beta) = \sum_{i=1}^N (y_i - p_i) x_i, 
\quad
H(\beta) = -\sum_{i=1}^N p_i(1-p_i) x_i x_i^T.
\]
Newton--Raphson update:
\[
\beta^{(t+1)} = \beta^{(t)} + \Big[ \sum_{i=1}^N p_i(1-p_i) x_i x_i^T \Big]^{-1} 
\sum_{i=1}^N (y_i - p_i)x_i.
\]
\begin{itemize}
    \item This can be written as a \textbf{weighted least squares problem}:  
    \[
    \beta^{(t+1)} = (X^T W X)^{-1} X^T W z,
    \]
    where:
    \[
    W = \mathrm{diag}(p_i(1-p_i)), \quad 
    z = X\beta^{(t)} + W^{-1}(y - p).
    \]
    \item Known as \textbf{Iteratively Reweighted Least Squares (IRLS)}.  
\end{itemize}
\end{frame}


%--------------------------------------------------
\begin{frame}{MLE: Second-Order Conditions}
Hessian (matrix of second derivatives):
\[
H(\beta) = \frac{\partial^2 \ell}{\partial \beta \partial \beta^T} 
= -\sum_{i=1}^N p_i(1-p_i) x_i x_i^T.
\]
\begin{itemize}
    \item Negative definite $\Rightarrow$ log-likelihood is concave.  
    \item Ensures unique global maximum.  
    \item Newton–Raphson update:
    \[
    \beta^{(t+1)} = \beta^{(t)} - H(\beta^{(t)})^{-1} \nabla \ell(\beta^{(t)}).
    \]
\end{itemize}
\end{frame}

%--------------------------------------------------
\begin{frame}{Asymptotic Properties of MLE}
Under standard regularity conditions:
\[
\hat{\beta} \xrightarrow{p} \beta_0 \quad \text{(consistency)}.
\]
\[
\sqrt{N}(\hat{\beta} - \beta_0) \xrightarrow{d} \mathcal{N}(0, I(\beta_0)^{-1}),
\]
where
\[
I(\beta_0) = \mathbb{E}\!\left[ p(X)(1-p(X))XX^T \right]
\]
is the Fisher information matrix.
\begin{itemize}
    \item Provides standard errors for inference.  
    \item Enables Wald tests, score tests, likelihood ratio tests.  
\end{itemize}
\end{frame}

%--------------------------------------------------
\begin{frame}{Summary of Logistic Regression}
\begin{itemize}
    \item Models $p(X) = \Pr(Y=1|X)$ using the logistic function.  
    \item Log-odds are linear in predictors.  
    \item Parameters estimated via maximum likelihood.  
    \item Gradient and Hessian yield efficient algorithms (Newton–Raphson, IRLS).  
    \item MLE has desirable large-sample properties (consistency, asymptotic normality).  
    \item Foundation for extensions: multinomial logit, generalized linear models.  
\end{itemize}
\end{frame}
\end{document}
